% Generated draft: requires mathematical and source review.
\documentclass[11pt]{article}
\usepackage[margin=25mm]{geometry}
\usepackage{fontspec}
\setmainfont{Latin Modern Roman}
\usepackage{amsmath,amssymb,mathtools,unicode-math}
\setmathfont{Latin Modern Math}
\usepackage{enumitem,needspace,xurl,hyperref}
\setlength{\parindent}{0pt}
\setlength{\parskip}{5pt}
\setlength{\emergencystretch}{3em}
\newcommand{\N}{\mathbb{N}}
\newcommand{\Z}{\mathbb{Z}}
\newcommand{\Q}{\mathbb{Q}}
\newcommand{\R}{\mathbb{R}}
\newcommand{\C}{\mathbb{C}}
\newcommand{\F}{\mathbb{F}}
\newcommand{\A}{\mathbb{A}}
\newcommand{\PP}{\mathbb{P}}
\newcommand{\E}{\mathbb{E}}
\newcommand{\eps}{\varepsilon}
\newcommand{\dd}{\,\mathrm{d}}
\begin{document}
\section*{UnsolvedMath ID: 5600008}
\textbf{Source problem code:} catalogue entry
\subsection{Definitions and mathematical statement}
Let $M$ be a compact surface and $I=(-1,1)$. Let $f:M\times I\to\R^3$ be differentiable and write $f_t(x)=f(x,t)$. Suppose that each $f_t$ is an immersion and that the induced first fundamental form is independent of $t$:
\[
\langle D_xf_t(v),D_xf_t(w)\rangle
=\langle D_xf_0(v),D_xf_0(w)\rangle
\quad(x\in M,\ v,w\in T_xM,\ t\in I).
\]
The literal question is whether there must be Euclidean rigid motions $g(t)$ such that $f_t=g(t)\circ f_0$ for every $t\in I$. Equivalently, it asks for functions $Q:I\to O(3)$ and $a:I\to\R^3$ satisfying
\[
f(x,t)=Q(t)f(x,0)+a(t)\qquad(x\in M,\ t\in I).
\]
No regularity of $g(t)$ beyond being a family is specified in the question.

The source does not impose convexity, embedding rather than immersion, connectedness, or a boundary convention for $M$. It says differentiable, while the title and supplied progress summary use smooth; precise differentiability thresholds matter to rigidity. These hypotheses must not be supplied silently. The question concerns a finite family of isometric immersions, not merely an infinitesimal deformation.

The supplied summary describes rigidity for smooth compact strictly convex hypersurfaces and failures without convexity. That restricted theorem does not prove the literal unrestricted assertion. Conversely, the existence of failures in some category does not identify their regularity, topology, and immersion/embedding hypotheses in this record.
\subsection{Short English statement}
Must every differentiable family of isometric immersions of a compact surface in three-space arise from rigid motions? Preserve the unrestricted source question while checking the missing regularity, convexity, embedding, boundary, and connectedness hypotheses.
\subsection{Sources}
\begin{itemize}
\item Local UnsolvedMath registry, permanent ID 5600008, exact source code catalogue entry, ``Rigidity of Smooth Isometric Families of Compact Surfaces''; proposed by Shiing-Shen Chern, 1970. Registry status open, classification OPEN-TRIAGE.
\item Shiing-Shen Chern, \emph{Open Problems in Differential Geometry} (1970), Problem 8, PDF page 9, historical ICM proceedings scan, \url{https://people.math.harvard.edu/~hirolee/pdfs/2014-fall-230a-icm1970-chern-differential-geometry.pdf}. Supplied normalized pdf-text extraction accessed July 29, 2026. No PDF was accessed here.
\item The supplied summary attributes strict-convexity rigidity to classical work of Pogorelov and Kuznetsov, and says the general nonconvex/nonpositive-curvature case admits failures. It provides no precise theorem or counterexample citation and calls the problem partially/largely resolved. These assertions are retained as qualified supplied assessments rather than an unqualified rigidity theorem.
\item Historical-source current status and statement-reuse terms for the public ICM scan are both flagged NEEDS\_REVIEW in the provenance.
\end{itemize}
\end{document}
